% =====================================================================
%  E9: E8's DB head set (J-head decoding of the task's most active words,
%  and little attention on BOS) on the five complex tasks, both checks,
%  alpha 1 and 2. Record of the experiment: why, how DB was carried over,
%  lists and sets, both checks, answer checks.
%  Tables are generated by scripts/e9_report.py (paper/tables/e9_*.tex).
% =====================================================================
\documentclass[10pt,a4paper]{article}

\usepackage[utf8]{inputenc}
\usepackage[T1]{fontenc}
\usepackage[margin=1.8cm]{geometry}
\usepackage{amsmath,amssymb}
\usepackage{booktabs,tabularx,array}
\usepackage{enumitem}
\usepackage{xcolor}
\usepackage{caption}
\usepackage{float}
\usepackage{graphicx}
\usepackage[colorlinks=true,linkcolor=blue!50!black,urlcolor=blue!60!black]{hyperref}

\setlist{nosep,leftmargin=*}
\setlength{\parindent}{0pt}
\setlength{\parskip}{5pt}
\renewcommand{\arraystretch}{1.12}
\captionsetup{font=small,labelfont=bf}

\title{E9: The DB Head Set on the Complex Tasks\\\large Lexical task heads, complex tasks}
\author{Filippo Merlo}
\date{10 November 2026}

\begin{document}
\maketitle
\vspace{-1.5em}

\paragraph{In short.} E9 carries E8's DB rule over to Fu et al.'s five complex tasks unchanged (heads that decode the task's 30 most active words under J-head and do not look mainly at BOS) and patches it in E2's two checks at $\alpha$ 1 and 2. DB holds 87--109 heads per task, more than E2's J-head sets, and most of its heads are also in another task's DB set. Like E2's sets, it shows no task-specific effect. Its one new behaviour is on the two Yes/No tasks. Without task information it makes the model give an answer at all, which no E2 set did, but nearly always ``Yes'', and right only as often as that label is. In the second check its gains there come from one label for every item, and the control mean gains the same. On Set Game and Tapatan the control mean matches or beats it.

% ---------------------------------------------------------------------
\section{Why}

On the simple tasks, E8's DB set (heads that decode the task's most active words under J-head and do not look mainly at BOS) was the best set that needs no hand-written word list and no causal step. On the complex tasks of Fu et al., E2's lexical task heads showed no task-specific effect in either check: the first check changed nothing on four tasks, and in the second the control mean matched every gain. E9 asks whether the DB rule, carried over unchanged, picks heads that do better there.

% ---------------------------------------------------------------------
\section{Design}

DB is built as in E8 and patched as in E2's complex checks. The changes to either are the ones the complex tasks force; they are listed in the last item.

\begin{enumerate}
\item \textbf{Prompts for the filters}: per task, the identification prompts of the seven rule wordings (original and w1--w6), pooled: each wording's first 100 prompts answered correctly under the flexible score, 700 prompts per task. Easy level, chat template, last prompt token, as the stored J-head readings of E2.
\item \textbf{D} decoding: share of these prompts whose J-head top 10 holds a token of the task's word list; pass if $\geq$0.10. List: the task's 30 most active decoded words under E8's rules (top-10 tokens decoded, lowercased, alphabetic, $\geq$3 letters, no stop words, Zipf $\geq$3; activity = share of (prompt, head) pairs whose top 10 holds the word), with the leading-space capitalised and upper-case forms, nothing removed.
\item \textbf{B} BOS: mean attention of the last prompt token on BOS over the same prompts $\leq$0.5. Measured in a new forward pass (eager attention, each prompt alone).
\item \textbf{Set}: DB = the heads that pass both. All layers, no size cap. Only DB is patched.
\item \textbf{Patching} (E2's complex checks): rule style; each rule wording's own mean (its mean set, 28--100 prompts) and the control task's mean under its original prompt (Noun Class Agreement and Lexical Category Inference control each other, Set Game and Operator Function too, Tapatan takes Set Game); $\alpha\in\{1,2\}$; replies to EOS or 1{,}024 tokens; flexible score; every answer saved. Check~1: the prompts without task information (the input block only) of the first 100 items answered correctly under the original rules prompt. Check~2: the original rules prompts of the first 100 items answered wrongly under it; Operator Function has none. Task accuracy = mean over the seven mean wordings.
\item \textbf{References} on the same items, means and replies: E2's rule-style J-head and Logit-head sets (Yang et al.'s rule with the hand-checked term lists, from the most accurate rule wording), at $\alpha$ 1 and 2 of their 11-value run.
\item \textbf{What differs from E8}: no fold split, since the complex checks have none, so both $\alpha$ are reported instead of a cross-fitted one; seven rule wordings pooled instead of five; prompts under the chat template, so BOS is followed by the template's header tokens; a control mean in both checks, as in E2. The word rules keep only English alphabetic words: digits (the answers of Operator Function), board symbols and the invented words of Noun Class Agreement cannot be on the lists.
\end{enumerate}

% ---------------------------------------------------------------------
\section{Word lists and sets}

\begin{table}[H]\centering\scriptsize
% Generated by scripts/e9_report.py; do not edit.
\begin{tabular}{@{}p{2.6cm}cp{12.4cm}@{}}
\toprule
Task & Generic & Words (most active first; generic in italics) \\
\midrule
Noun Class Agreement & 14 & yes, answers, \textit{http}, answer, nos, \textit{what}, answered, grammar, non, linguistic, \textit{why}, \textit{only}, gram, \textit{amongst}, \textit{because}, verdict, neither, \textit{mph}, covid, wikipedia, \textit{however}, \textit{which}, \textit{our}, \textit{etc}, answering, \textit{https}, nothing, \textit{their}, nazis, \textit{then} \\
Lexical Category Inference & 13 & yes, answer, category, \textit{http}, \textit{what}, list, referencing, \textit{mph}, \textit{https}, likely, \textit{which}, covid, answers, nos, \textit{tom}, wikipedia, \textit{only}, \textit{why}, \textit{because}, lists, \textit{amongst}, \textit{our}, categories, name, answered, today, according, \textit{additionally}, puzzle, \textit{etc} \\
Set Game & 5 & card, cards, puzzle, \textit{https}, first, three, based, game, \textit{http}, third, according, pairing, \textit{which}, pairs, puzzles, myself, covid, board, here, likely, games, analysis, pair, transitioning, second, each, \textit{respectively}, two, trio, \textit{only} \\
Operator Function & 8 & problem, rules, math, rule, according, mathematical, therefore, first, covid, recipe, \textit{what}, calculation, \textit{our}, using, problems, \textit{mph}, \textit{http}, \textit{amongst}, \textit{which}, calculations, \textit{additionally}, example, calculator, procedure, mathematics, wikipedia, \textit{https}, formula, aforementioned, applying \\
Tapatan & 6 & \textit{http}, player, wikipedia, \textit{https}, \textit{mph}, continue, \textit{what}, likely, board, covid, based, players, reddit, games, \textit{there}, final, \textit{however}, yesterday, horizontal, myself, algorithm, game, therefore, current, according, now, theater, when, biography, analysis \\
\bottomrule
\end{tabular}

\caption{D's word lists, most active first. Generic (italics): among the 50 words heads decode most on WikiText under J-head (E8's definition).}
\label{tab:lists}
\end{table}

As on the simple tasks, a part of each list is generic (5--14 of 30 words in the WikiText sense; \emph{covid}, \emph{wikipedia} or \emph{reddit} look generic too but are not in that set). The rest names the answer format (\emph{yes}, \emph{answer}, \emph{nos} on the two Yes/No tasks), the task's objects (\emph{card}, \emph{board}, \emph{player}, \emph{continue}) or its kind (\emph{math}, \emph{calculation}, \emph{grammar}, \emph{category}). No list holds the items a task is applied to (Set Game's animals and habitats, the words of the synthetic language, the objects Lexical Category Inference sorts); words like \emph{card}, \emph{board} or \emph{category} are in the prompt itself. Operator Function's answers are numbers, which the word rules leave out.

\begin{table}[H]\centering\small
\resizebox{\linewidth}{!}{% Generated by scripts/e9_report.py; do not edit.
\begin{tabular}{@{}lrrrrcccrr@{}}
\toprule
Task & Prompts & D & B & DB & Layers (median) & in E2 J & in E2 L & in another DB & in all DB \\
\midrule
Noun Class Agreement & 700 & 190 & 444 & 99 & 6--31 (16) & 8/26 & 2/15 & 85 & 14 \\
Lexical Category Inference & 700 & 192 & 449 & 109 & 3--31 (16) & 21/53 & 13/43 & 99 & 14 \\
Set Game & 700 & 184 & 460 & 97 & 3--31 (17) & 33/75 & 20/55 & 77 & 14 \\
Operator Function & 700 & 195 & 431 & 87 & 3--31 (15) & 9/40 & 5/23 & 67 & 14 \\
Tapatan & 700 & 184 & 450 & 105 & 1--31 (16) & 25/77 & 20/61 & 87 & 14 \\
\bottomrule
\end{tabular}
}
\caption{Heads per filter and in DB (1{,}024 heads in all); layers of the DB heads. In E2 J / L: DB heads in E2's rule-style J-head / Logit-head set, out of that set's size. In another DB: DB heads also in the DB set of at least one other task; in all DB: heads in all five DB sets.}
\label{tab:sets}
\end{table}

D passes about a fifth of all heads and B about 45\%, so DB keeps roughly half of D. The sets spread over all layers, with a median of 15--17. They are larger than E2's J-head sets (1.3 to 3.8 times) and hold a fifth to a little under half of them. They are not specific to a task: on every task, more than three quarters of the DB heads are also in another task's DB set, and 14 heads are in all five.

% ---------------------------------------------------------------------
\section{Results}

\subsection{Check 1: prompts without task information}

\begin{table}[H]\centering\small
% Generated by scripts/e9_report.py; do not edit.
\begin{tabular}{@{}lrccl r cc cc@{}}
\toprule
 & & & & & & \multicolumn{2}{c}{Task mean} & \multicolumn{2}{c}{Control mean} \\
\cmidrule(lr){7-8}\cmidrule(lr){9-10}
Task & Items & Unp. & Const. & Set & Heads & $\alpha{=}1$ & $\alpha{=}2$ & $\alpha{=}1$ & $\alpha{=}2$ \\
\midrule
Noun Class Agreement & 100 & 0.00 & 0.57 & DB (E9) & 99 & 0.06 & 0.21 & 0.02 & 0.00 \\
 &  &  &  & J-head (E2) & 26 & 0.00 & 0.00 & 0.00 & 0.00 \\
 &  &  &  & Logit-head (E2) & 15 & 0.00 & 0.00 & 0.00 & 0.00 \\
\addlinespace
Lexical Category Inference & 100 & 0.00 & 0.59 & DB (E9) & 109 & 0.03 & 0.13 & 0.32 & 0.41 \\
 &  &  &  & J-head (E2) & 53 & 0.00 & 0.00 & 0.00 & 0.00 \\
 &  &  &  & Logit-head (E2) & 43 & 0.00 & 0.00 & 0.00 & 0.00 \\
\addlinespace
Set Game & 100 & 0.58 & -- & DB (E9) & 97 & 0.59 & 0.57 & 0.62 & 0.39 \\
 &  &  &  & J-head (E2) & 75 & 0.60 & 0.58 & 0.63 & 0.58 \\
 &  &  &  & Logit-head (E2) & 55 & 0.60 & 0.59 & 0.55 & 0.61 \\
\addlinespace
Operator Function & 100 & 0.00 & -- & DB (E9) & 87 & 0.00 & 0.00 & 0.00 & 0.00 \\
 &  &  &  & J-head (E2) & 40 & 0.00 & 0.00 & 0.00 & 0.00 \\
 &  &  &  & Logit-head (E2) & 23 & 0.00 & 0.00 & 0.00 & 0.00 \\
\addlinespace
Tapatan & 100 & 0.00 & 0.58 & DB (E9) & 105 & 0.01 & 0.01 & 0.01 & 0.01 \\
 &  &  &  & J-head (E2) & 77 & 0.00 & 0.00 & 0.00 & 0.00 \\
 &  &  &  & Logit-head (E2) & 61 & 0.00 & 0.01 & 0.00 & 0.00 \\
\bottomrule
\end{tabular}

\caption{Check 1. Accuracy of the patched replies, flexible score, task mean = mean over the seven rule wordings' means. Items: patched items; Unp.: their unpatched accuracy; Const.: accuracy of answering every item with its most frequent label (label tasks). DB (E9) and E2's rule-style J-head and Logit-head sets on the same items and means, at $\alpha$ 1 and 2. Control mean: the control task's mean under its original prompt, at the same heads.}
\label{tab:check1}
\end{table}

\begin{itemize}
\item \textbf{Operator Function, Tapatan}: nothing changes. As in E2, the model asks for the function (Operator Function) or describes the board without a label (Tapatan), under every mean.
\item \textbf{Set Game}: the task means leave accuracy where it is without a patch. The control mean at $\alpha=2$ lowers it, with 40\% of its replies giving no readable set.
\item \textbf{Noun Class Agreement}: DB is the first set that gets an answer out of these prompts: with the task means, the model answers 13\% of the replies at $\alpha=1$ and 45\% at $\alpha=2$, always with ``Yes'' (Table~\ref{tab:labels}). It is right exactly on those of the answered items whose reference is ``Yes''. These are about half of the answered items, close to their share among all items (43\%), so which prompts get an answer hardly depends on the right label. The task means bring back the answer format with one label, not the task. The control mean does almost nothing here.
\item \textbf{Lexical Category Inference}: the task means mostly leave the prompts unanswered, except the means of two wordings (w1, w4), which at $\alpha=2$ answer ``Yes'' to nearly every item (Appendix, Table~\ref{tab:wordings}). The control mean, Noun Class Agreement's, does more: at $\alpha=2$ it answers ``Yes'' to every item, which scores the share of ``Yes''. Only at $\alpha=1$, where it answers half the prompts with both labels, is it right somewhat more often than one label would be on the same items.
\end{itemize}

\begin{table}[H]\centering\scriptsize
\resizebox{\linewidth}{!}{% Generated by scripts/e9_report.py; do not edit.
\begin{tabular}{@{}llll cccc cccc@{}}
\toprule
 & & & & \multicolumn{4}{c}{$\alpha{=}1$} & \multicolumn{4}{c}{$\alpha{=}2$} \\
\cmidrule(lr){5-8}\cmidrule(lr){9-12}
Check & Task & Set & Mean & Answered & Top answer & Right & One label & Answered & Top answer & Right & One label \\
\midrule
check 1 & Noun Class Agreement & DB (E9) & task & 0.13 & Yes 0.13 & 0.06 & 0.06 & 0.45 & Yes 0.45 & 0.21 & 0.24 \\
 &  &  & control & 0.05 & Yes 0.05 & 0.02 & 0.03 & 0.00 & -- & 0.00 & 0.00 \\
 &  & J-head (E2) & task & 0.00 & -- & 0.00 & 0.00 & 0.00 & -- & 0.00 & 0.00 \\
 &  &  & control & 0.00 & -- & 0.00 & 0.00 & 0.00 & -- & 0.00 & 0.00 \\
\addlinespace
 & Lexical Category Inference & DB (E9) & task & 0.05 & No 0.04 & 0.03 & 0.03 & 0.28 & Yes 0.28 & 0.13 & 0.17 \\
 &  &  & control & 0.51 & Yes 0.37 & 0.32 & 0.27 & 1.00 & Yes 1.00 & 0.41 & 0.59 \\
 &  & J-head (E2) & task & 0.00 & -- & 0.00 & 0.00 & 0.00 & -- & 0.00 & 0.00 \\
 &  &  & control & 0.00 & -- & 0.00 & 0.00 & 0.00 & -- & 0.00 & 0.00 \\
\addlinespace
 & Tapatan & DB (E9) & task & 0.01 & continue 0.01 & 0.01 & 0.01 & 0.01 & continue 0.01 & 0.01 & 0.01 \\
 &  &  & control & 0.01 & continue 0.01 & 0.01 & 0.01 & 0.02 & continue 0.02 & 0.01 & 0.01 \\
 &  & J-head (E2) & task & 0.01 & continue 0.01 & 0.00 & 0.00 & 0.00 & continue 0.00 & 0.00 & 0.00 \\
 &  &  & control & 0.00 & -- & 0.00 & 0.00 & 0.00 & -- & 0.00 & 0.00 \\
\midrule
check 2 & Noun Class Agreement & DB (E9) & task & 1.00 & No 0.94 & 0.36 & 0.58 & 1.00 & Yes 1.00 & 0.58 & 0.58 \\
 &  &  & control & 1.00 & No 0.93 & 0.35 & 0.58 & 1.00 & Yes 1.00 & 0.58 & 0.58 \\
 &  & J-head (E2) & task & 1.00 & No 0.59 & 0.03 & 0.58 & 1.00 & No 0.63 & 0.05 & 0.58 \\
 &  &  & control & 1.00 & No 0.67 & 0.09 & 0.58 & 1.00 & No 0.74 & 0.16 & 0.58 \\
\addlinespace
 & Lexical Category Inference & DB (E9) & task & 1.00 & No 0.81 & 0.33 & 0.72 & 1.00 & No 0.71 & 0.41 & 0.72 \\
 &  &  & control & 1.00 & No 1.00 & 0.28 & 0.72 & 1.00 & No 0.84 & 0.38 & 0.72 \\
 &  & J-head (E2) & task & 1.00 & No 0.77 & 0.12 & 0.72 & 1.00 & No 0.99 & 0.27 & 0.72 \\
 &  &  & control & 1.00 & No 0.84 & 0.12 & 0.72 & 1.00 & No 0.99 & 0.27 & 0.72 \\
\addlinespace
 & Tapatan & DB (E9) & task & 0.84 & A win 0.45 & 0.17 & 0.36 & 0.88 & A win 0.56 & 0.19 & 0.38 \\
 &  &  & control & 0.85 & A win 0.53 & 0.17 & 0.35 & 0.88 & A win 0.46 & 0.31 & 0.39 \\
 &  & J-head (E2) & task & 0.86 & A win 0.52 & 0.15 & 0.37 & 0.85 & A win 0.53 & 0.16 & 0.35 \\
 &  &  & control & 0.86 & A win 0.55 & 0.16 & 0.38 & 0.90 & A win 0.54 & 0.15 & 0.39 \\
\bottomrule
\end{tabular}
}
\caption{Answers on the label tasks, pooled over the means of each condition (task: 7 wordings $\times$ 100 replies; control: 100). Answered: share of replies with an answer under the flexible score; top answer: the most frequent one, with its share of all replies; right: accuracy; one label: the accuracy of giving every answered item the label most frequent among the references of the answered items (unanswered items count as wrong), the line a one-label shift reaches without the task.}
\label{tab:labels}
\end{table}

\subsection{Check 2: original prompts answered wrongly}

\begin{table}[H]\centering\small
% Generated by scripts/e9_report.py; do not edit.
\begin{tabular}{@{}lrccl r cc cc@{}}
\toprule
 & & & & & & \multicolumn{2}{c}{Task mean} & \multicolumn{2}{c}{Control mean} \\
\cmidrule(lr){7-8}\cmidrule(lr){9-10}
Task & Items & Unp. & Const. & Set & Heads & $\alpha{=}1$ & $\alpha{=}2$ & $\alpha{=}1$ & $\alpha{=}2$ \\
\midrule
Noun Class Agreement & 100 & 0.00 & 0.58 & DB (E9) & 99 & 0.36 & 0.58 & 0.35 & 0.58 \\
 &  &  &  & J-head (E2) & 26 & 0.03 & 0.05 & 0.09 & 0.16 \\
 &  &  &  & Logit-head (E2) & 15 & 0.04 & 0.05 & 0.08 & 0.13 \\
\addlinespace
Lexical Category Inference & 100 & 0.00 & 0.72 & DB (E9) & 109 & 0.33 & 0.41 & 0.28 & 0.38 \\
 &  &  &  & J-head (E2) & 53 & 0.12 & 0.27 & 0.12 & 0.27 \\
 &  &  &  & Logit-head (E2) & 43 & 0.09 & 0.28 & 0.04 & 0.21 \\
\addlinespace
Set Game & 100 & 0.00 & -- & DB (E9) & 97 & 0.21 & 0.21 & 0.20 & 0.24 \\
 &  &  &  & J-head (E2) & 75 & 0.21 & 0.22 & 0.18 & 0.23 \\
 &  &  &  & Logit-head (E2) & 55 & 0.22 & 0.22 & 0.24 & 0.23 \\
\addlinespace
Operator Function & \multicolumn{9}{l}{\textit{no item answered wrongly under orig}} \\
Tapatan & 100 & 0.00 & 0.43 & DB (E9) & 105 & 0.17 & 0.19 & 0.17 & 0.31 \\
 &  &  &  & J-head (E2) & 77 & 0.15 & 0.16 & 0.16 & 0.15 \\
 &  &  &  & Logit-head (E2) & 61 & 0.13 & 0.13 & 0.13 & 0.17 \\
\bottomrule
\end{tabular}

\caption{Check 2: the original rules prompts of the first 100 items answered wrongly under them (unpatched accuracy 0 by construction), as Table~\ref{tab:check1}.}
\label{tab:check2}
\end{table}

\begin{itemize}
\item \textbf{Noun Class Agreement}: DB's gains are far above E2's, and the control mean's are the same. At $\alpha=2$ every mean, task or control, answers ``Yes'' to (nearly) every item, which is the constant-answer line; at $\alpha=1$ most task means answer ``No'' to every item (Table~\ref{tab:wordings}). The task means and the control are right on nearly the same items (Table~\ref{tab:overlap}).
\item \textbf{Lexical Category Inference}: per wording, accuracy is at or near one of two values, the share of ``No'' or the share of ``Yes'' among the items (Table~\ref{tab:wordings}): each mean gives one label to almost every item. The mean over wordings is a little above the control because one or two of the seven wordings' means push towards ``Yes'', the more frequent label here. E2's sets push towards ``No'' and score less.
\item \textbf{Set Game}: DB, E2's sets and the control mean are level. The task means and the control are right on different items, with about half shared, as expected when a patch sends the reasoning down another path rather than adding task information (proposal-complex, ``A control for the second check'').
\item \textbf{Tapatan}: the control mean at $\alpha=2$ is above every task mean; at $\alpha=1$ they are level.
\end{itemize}

\begin{table}[H]\centering\small
% Generated by scripts/e9_report.py; do not edit.
\begin{tabular}{@{}ll ccc ccc@{}}
\toprule
 & & \multicolumn{3}{c}{$\alpha{=}1$} & \multicolumn{3}{c}{$\alpha{=}2$} \\
\cmidrule(lr){3-5}\cmidrule(lr){6-8}
Task & Set & Task mean & Control & Both & Task mean & Control & Both \\
\midrule
Noun Class Agreement & DB (E9) & 36.0 & 35 & 31.1 & 57.7 & 58 & 57.7 \\
 & J-head (E2) & 3.1 & 9 & 2.0 & 5.4 & 16 & 5.4 \\
Lexical Category Inference & DB (E9) & 32.9 & 28 & 21.1 & 40.6 & 38 & 21.6 \\
 & J-head (E2) & 11.9 & 12 & 7.0 & 27.4 & 27 & 26.7 \\
Set Game & DB (E9) & 20.6 & 20 & 9.4 & 21.1 & 24 & 10.1 \\
 & J-head (E2) & 21.0 & 18 & 10.1 & 22.4 & 23 & 13.6 \\
Tapatan & DB (E9) & 16.7 & 17 & 5.1 & 18.7 & 31 & 10.3 \\
 & J-head (E2) & 14.6 & 16 & 9.6 & 16.4 & 15 & 8.4 \\
\bottomrule
\end{tabular}

\caption{Check 2: items right under a task mean (mean over the seven wordings), under the control mean, and under both (mean over wordings), out of 100.}
\label{tab:overlap}
\end{table}

% ---------------------------------------------------------------------
\section{Takeaways}

\begin{enumerate}
\item Carried over unchanged, the DB rule finds no task-specific heads on the complex tasks: no gain of a task mean in either check that the control mean or a one-label shift does not explain. E2's conclusion holds for this selection too.
\item On the complex prompts the decoded words that pick the heads are mostly answer-format words, the task's objects and generic words, and the resulting sets are largely shared across tasks.
\item What DB adds over E2's sets is a stronger push towards one label on the Yes/No tasks: in check~1 it brings the answer format back (with ``Yes''), in check~2 it gives every item one label. In check~2 the control mean does the same at DB's heads; in check~1 the push comes with Noun Class Agreement's means, on its own prompts and, as the control, on Lexical Category Inference's. So it is a property of DB's heads and of those means, not of each task's own information. DB is also larger than E2's sets on these tasks, so size may be part of it.
\end{enumerate}

\paragraph{Caveats.} One run per setting, greedy, no intervals, no random or size-matched head sets. Only $\alpha$ 1 and 2: E2's largest effects (one-label shifts, Set Game at large $\alpha$) came at $|\alpha|$ of 4 to 10, so this run can say nothing about them. No fold split, as in E2: the heads are chosen on the rule prompts of the same items whose bare prompts are patched in check~1, and the means of the six new wordings contain 40--89 of the 100 items patched in check~2 (prompts answered rightly under that wording, wrongly under the original). The word rules leave out digits and the invented words of the synthetic language, and the BOS threshold comes from the simple tasks; both were kept, not tuned. Set sizes differ from E2's.

\paragraph{Open next steps.}
\begin{itemize}
\item Random head sets of DB's size and layer profile, to see whether the stronger one-label shift on the Yes/No tasks comes from set size.
\item The 14 heads common to all five DB sets alone, in check~2 of the Yes/No tasks.
\end{itemize}

\paragraph{Reproduce.}
\begin{itemize}
\item \texttt{scripts/e9\_attention.py}: the BOS masses (GPU, about 30 s per task).
\item \texttt{scripts/e9\_sets.py --out results/e9\_complex\_db}: word lists and sets (no GPU).
\item \texttt{scripts/e9\_patch.sh}: both checks (\texttt{scripts/patch.py --set-list --alphas 1 2}; GPU, about 1 h on 4 GPUs).
\item \texttt{scripts/e9\_report.py}: accuracies, answer checks, \texttt{results/e9\_complex\_db/report.md} and these tables.
\end{itemize}

\appendix
\section{Per wording}

\begin{table}[H]\centering\small
% Generated by scripts/e9_report.py; do not edit.
\begin{tabular}{@{}llcccccccc@{}}
\toprule
Check & Task & $\alpha$ & orig & w1 & w2 & w3 & w4 & w5 & w6 \\
\midrule
check 1 & Noun Class Agreement & 1 & 0.05 & 0.04 & 0.08 & 0.09 & 0.09 & 0.04 & 0.05 \\
 &  & 2 & 0.13 & 0.16 & 0.20 & 0.31 & 0.36 & 0.15 & 0.17 \\
 & Lexical Category Inference & 1 & 0.00 & 0.09 & 0.00 & 0.00 & 0.15 & 0.00 & 0.00 \\
 &  & 2 & 0.01 & 0.44 & 0.03 & 0.00 & 0.41 & 0.00 & 0.00 \\
 & Set Game & 1 & 0.55 & 0.61 & 0.56 & 0.62 & 0.59 & 0.62 & 0.61 \\
 &  & 2 & 0.58 & 0.60 & 0.59 & 0.54 & 0.60 & 0.55 & 0.55 \\
 & Operator Function & 1 & 0.00 & 0.00 & 0.00 & 0.00 & 0.00 & 0.00 & 0.00 \\
 &  & 2 & 0.00 & 0.00 & 0.00 & 0.00 & 0.00 & 0.00 & 0.00 \\
 & Tapatan & 1 & 0.01 & 0.00 & 0.01 & 0.01 & 0.01 & 0.01 & 0.01 \\
 &  & 2 & 0.01 & 0.01 & 0.01 & 0.01 & 0.01 & 0.00 & 0.01 \\
\midrule
check 2 & Noun Class Agreement & 1 & 0.21 & 0.42 & 0.41 & 0.42 & 0.42 & 0.22 & 0.42 \\
 &  & 2 & 0.58 & 0.58 & 0.58 & 0.58 & 0.58 & 0.58 & 0.56 \\
 & Lexical Category Inference & 1 & 0.28 & 0.28 & 0.21 & 0.28 & 0.27 & 0.70 & 0.28 \\
 &  & 2 & 0.28 & 0.28 & 0.72 & 0.28 & 0.28 & 0.72 & 0.28 \\
 & Set Game & 1 & 0.16 & 0.22 & 0.20 & 0.19 & 0.23 & 0.22 & 0.22 \\
 &  & 2 & 0.28 & 0.18 & 0.18 & 0.21 & 0.24 & 0.22 & 0.17 \\
 & Tapatan & 1 & 0.16 & 0.16 & 0.16 & 0.18 & 0.21 & 0.15 & 0.15 \\
 &  & 2 & 0.18 & 0.14 & 0.21 & 0.22 & 0.19 & 0.17 & 0.20 \\
\bottomrule
\end{tabular}

\caption{DB, task accuracy under each rule wording's mean.}
\label{tab:wordings}
\end{table}

\end{document}
